% ═══════════════════════════════════════════════════════════════════
%  Results
% ═══════════════════════════════════════════════════════════════════

\label{sec:results}

This section reports results for Experiments~1--6 (including the Experiment~4 anchor-environment extension) using \texttt{Qwen/Qwen2.5-14B-Instruct} with the free-language agent type and the feasibility gate disabled throughout.  Section~\ref{sec:results-quality} validates the experimental backend before any behavioural claims are made.  Section~\ref{sec:results-micro} reports micro-level findings for individual bargaining behaviour.  Section~\ref{sec:results-macro} reports macro-level findings for market outcomes.  Section~\ref{sec:results-cross} synthesises cross-experiment patterns, deal success rates, and behavioural comparisons with an explicit evidence-strength assessment.

% ═══════════════════════════════════════════════════════════════════
%  Pipeline Validity and Sanity Checks
% ═══════════════════════════════════════════════════════════════════

\subsection{Pipeline Validity and Sanity Checks}
\label{sec:results-quality}

\paragraph{Backend validity.}
The \texttt{Qwen2.5-14B-Instruct} model produced valid, parseable responses across all sessions (parse success rate: 100\%).  The feasibility gate is disabled throughout, so every acceptance is LLM-driven rather than gate-overridden.  Across the core six experiments (1--6; 2,220 nominal sessions, 2,171 after surplus-max matcher exclusions), the overall deal rate is 78.6\% (1,705 deals); including the Experiment~4 anchor-environment extension (1,200 additional sessions), the total programme comprises 3,420 sessions.  Non-deal outcomes are either timeouts (maximum rounds reached without agreement) or implicit rejections through counter-offer exhaustion.

Because the gate is disabled throughout, results are interpretable as agent behaviour rather than as the combined output of agent decisions and deterministic enforcement.  The counteroffer protocol, acceptance hints, and monotonic concession instructions remain active as fixed protocol parameters across all conditions (Section~\ref{sec:confounds}).

% ═══════════════════════════════════════════════════════════════════
%  Micro-Level Results
% ═══════════════════════════════════════════════════════════════════

\subsection{Micro-Level Results: Agent Bargaining Behaviour}
\label{sec:results-micro}

Experiments~1 and~2 use fixed-scenario parameters to isolate individual agent behaviour from market-level variation.

% ───────────────────────────────────────────────────────────────────
\subsubsection{Experiment 1: Concession Dynamics}
\label{sec:results-concession}

Table~\ref{tab:concession} reports the full concession trajectory for the rule-based and \texttt{llm\_free\_language} conditions.  The scenario is fixed at buyer value~=~\$120, seller cost~=~\$80, giving a ZOPA of \$40.  All 140 deals across 144 sessions are LLM-driven acceptances; zero gate-driven settlements occur.  The deal rate is 97.2\% (140/144).

\begin{table}[ht]
\centering
\caption{Concession trajectories by condition and round ($n = 72$ per condition, 3 seeds).  LLM values are mean $\pm$ standard deviation; rule-based values are deterministic (zero variance).  Buyer value~=~\$120, seller cost~=~\$80.}
\label{tab:concession}
\begin{tabular}{@{}lcccccc@{}}
\toprule
 & \multicolumn{3}{c}{Buyer offers (\$)} & \multicolumn{3}{c}{Seller counters (\$)} \\
\cmidrule(lr){2-4}\cmidrule(lr){5-7}
Condition & Round 0 & Round 2 & Round 4 & Round 1 & Round 3 & Round 5 \\
\midrule
\texttt{rule\_based}       & 60.00 & 73.33 & 86.67 & 101.33 & 96.00 & 90.67 \\
\texttt{llm\_free\_language} & $71.88 \pm 3.79$ & $85.18 \pm 0.87$ & $86.72 \pm 1.93$ & $89.93 \pm 0.59$ & $87.60 \pm 2.05$ & $88.06 \pm 1.87$ \\
\bottomrule
\end{tabular}
\end{table}

\paragraph{Finding 1: LLM buyer concession follows a concave trajectory.}
The LLM free-language buyer opens at \$71.88 (round~0), jumps by \$13.31 to \$85.18 in round~2, then slows sharply: round~4 moves by only \$1.54, with subsequent steps of \$0.31 and \$0.14.  This concave deceleration---large initial concession followed by rapid hardening---is consistent with Boulware-style bargaining behaviour, where an agent makes early accommodations but becomes increasingly resistant as the negotiation progresses.  The pattern is theoretically coherent: a large initial concession signals genuine willingness to negotiate, while subsequent hardening communicates commitment to a position---a ``tough but not foolish'' strategy that preserves credibility.  The magnitude of the initial jump is notable: at \$13.31, it represents approximately one-third of the full ZOPA width (\$40) in a single move, suggesting that the agent front-loads its flexibility early rather than distributing it evenly.  The rule-based buyer, by contrast, concedes at a fixed linear rate of approximately \$13.33 per round.  H-1b is partially supported: LLM agents settle earlier than the rule-based agent, though not because of a more even concession schedule.

\paragraph{Finding 2: Concession burden is asymmetrically borne by the buyer.}
Mean deal price is \$87.82.  Against the ZOPA midpoint of \$100 (equidistant between buyer value \$120 and seller cost \$80), the buyer captures \$32.18 in surplus (80.5\% of ZOPA) while the seller captures only \$7.82 (19.5\%).  The buyer-to-seller concession ratio is 5.65:1.  All deals cluster in the \$80--\$91.25 price range, occupying only the bottom 28\% of the available ZOPA.  This asymmetry reflects differential concession willingness: the LLM seller exhibits strong anchor rigidity, treating its early counter-offer as a defensible position and resisting further movement, while the buyer bears the concession burden to close the deal.  This pattern is behaviourally consistent with the anchoring literature~\cite{galinsky2001}, in which the party that establishes an early positional anchor captures a disproportionate share of the surplus.  A subtle but important feature of this result is that the effective anchor is the \emph{seller's} counter-offer rather than the buyer's opening bid, even though the buyer moves first.  The buyer's initial offer (\$71.88) serves as a starting point for movement, but the seller's counter at \$89.93 appears to function as the positional anchor around which the final deal is reached---the seller barely moves from it (\$2 total concession) while the buyer closes nearly the entire remaining gap.  This is consistent with Galinsky and Mussweiler's~\cite{galinsky2001} finding that the party whose reference point governs the negotiation captures more surplus, independent of who makes the first move: what matters is whose offer becomes the de facto anchor, not the turn order.

\paragraph{Finding 3: Seller anchors high and concedes slowly.}
The LLM seller opens at \$89.93 in round~1---notably lower than the rule-based seller's \$101.33---but then barely moves: rounds~3 and~5 are \$87.60 and \$88.06, a total seller concession of under \$2.  Seller variance is also low ($\pm$\$0.59 at round~1), indicating consistent behaviour across seeds.  The seller appears to find an early defensible position and resist further movement.

\paragraph{Rule-based baseline observation.}
The rule-based buyer opens at exactly 50\% of its value ceiling (\$60.00) and moves linearly, while the LLM buyer opens \$11.88 higher.  The LLM buyer's willingness to make a more generous opening offer is directionally consistent with H-1c (structured reasoning affects offer behaviour), though a direct comparison between rule-based and free-language conditions is confounded by the complete difference in decision architecture.

% ───────────────────────────────────────────────────────────────────
\subsubsection{Experiment 2: Deadline Pressure}
\label{sec:results-deadline}

Table~\ref{tab:deadline} reports settlement timing and price outcomes for three time-horizon conditions (maximum rounds: 6, 12, 20) with the gate disabled and \texttt{llm\_free\_language} agents.  216 sessions were run (72 per condition); deal rate is 96.3\% (208 deals).  All 208 deals are LLM-driven.

\begin{table}[ht]
\centering
\caption{Deadline experiment outcomes by time horizon ($n = 72$ per condition, 3 seeds).  Gate disabled; buyer value~=~\$120, seller cost~=~\$80.  Concession rate is mean \$/round to settlement.}
\label{tab:deadline}
\begin{tabular}{@{}lccccc@{}}
\toprule
Max rounds & Deal rate & Mean price (\$) & Avg rounds & Last-2-round share & Concession rate (\$/round) \\
\midrule
6  & 94.4\% & 86.47 & 5.21 & 82.4\% & 2.73 \\
12 & 97.2\% & 86.73 & 6.70 & 15.7\% & 2.36 \\
20 & 97.2\% & 84.93 & 7.30 & 11.4\% & 2.18 \\
\bottomrule
\end{tabular}
\end{table}

\paragraph{Finding 4: Deadline pressure accelerates concession.}
Under the 6-round horizon, 82.4\% of all deals close in the last two rounds, compared to 15.7\% under 12 rounds and 11.4\% under 20 rounds.  The tight-deadline condition also produces the fastest concession rate (2.73~\$/round vs.\ 2.18~\$/round under 20 rounds), a 25\% increase.  H-2a is supported.  The result is directionally consistent with deadline effects documented in human negotiation~\cite{roth1988deadline}.

\paragraph{Finding 5: Prices are stable across horizon conditions.}
Mean deal prices differ by less than \$2 across conditions (\$84.93 to \$86.73), indicating that deadline proximity accelerates settlement timing without materially changing the final price.  Agents concede faster but do not concede more: the surplus split remains approximately 83--88\% buyer-favourable across all three conditions.  This is consistent with the asymmetric concession pattern observed in Experiment~1: buyers absorb most of the concession pressure, but even under tight deadlines the seller price remains close to cost.  The price stability across deadline conditions constitutes a theoretically interesting null result.  In human deadline studies~\cite{roth1988deadline}, deadline pressure typically affects both the timing of concessions and the magnitude of final concessions, with tighter deadlines producing lower-quality outcomes for one or both parties.  The LLM result dissociates these two effects: timing is sensitive to the horizon (reflected in the last-two-round share shifting from 11.4\% to 82.4\%), but price is not.  A plausible interpretation is that the LLM agent treats deadline salience as a ``time to close'' signal---an instruction to accelerate agreement---rather than as a ``concede more'' signal that would alter its positional target.  This dissociation is a specific, testable behavioural difference between LLM and human negotiators that could be probed further through targeted deadline studies with human baselines.

% ═══════════════════════════════════════════════════════════════════
%  Macro-Level Results
% ═══════════════════════════════════════════════════════════════════

\subsection{Macro-Level Results: Market Outcomes and Mechanism Effects}
\label{sec:results-macro}

Experiments~3, 4, 6, the Experiment~4 extension, and~5 shift the unit of analysis from individual sessions to multi-tick markets.  All use \texttt{llm\_free\_language} agents with the gate disabled and 3 seeds.

% ───────────────────────────────────────────────────────────────────
\subsubsection{Experiment 3: Market Dynamics}
\label{sec:results-market}

Experiment~3 observes a single-condition 10-tick market (10 pairs per tick, 3 seeds = 300 sessions total) to examine whether bilateral LLM negotiation produces stable aggregate price levels and whether within-tick price dispersion persists.  The overall deal rate is 82.3\% (247/300); the timeout rate is 9.7\%.

\begin{table}[ht]
\centering
\caption{Market dynamics tick-level statistics averaged across seeds (300 sessions, 3 seeds).  CV is coefficient of variation of mean price across ticks.}
\label{tab:market-agg}
\begin{tabular}{@{}lcccc@{}}
\toprule
Tick & Mean price (\$) & Price std (\$) & Liquidity & Notes \\
\midrule
0 & 83.68 & 14.54 & --- & Early \\
1 & 97.29 & 17.53 & --- & \\
5 & 92.16 & 5.51  & --- & Mid \\
9 & 84.50 & 16.51 & --- & Late \\
\addlinespace
\multicolumn{5}{@{}l}{\textit{Aggregate (all 10 ticks)}} \\
Mean price (all deals) & \multicolumn{2}{l}{$\sim$\$89} & --- & --- \\
Mean liquidity & --- & --- & 0.81 & High \\
Price trend slope & \multicolumn{2}{l}{$-$0.60 per tick} & --- & Mild downward drift \\
\bottomrule
\end{tabular}
\end{table}

\paragraph{Finding 6: Prices do not converge; dispersion persists throughout.}
The coefficient of variation of tick-level mean prices fluctuates between 0.06 and 0.20 across the 10-tick window without a consistent downward trend.  The market does not settle on a common price level, consistent with the structural property of this market design: each pair negotiates independently under private constraints, and no mechanism forces prices to converge across pairs.  Within-tick price standard deviation ranges from \$5.51 to \$17.53, confirming that the Law of One Price fails in this market structure.  H-3a is not supported: mean price does not stabilise within the observed horizon.  H-3b is supported: within-tick dispersion remains substantial throughout.  Importantly, the absence of price convergence is not a failure of the market or the agents---it is the theoretically predicted outcome for a system of independent bilateral negotiations.  In such a structure, no price-discovery mechanism connects the outcomes of simultaneous pairs; each negotiation is resolved privately, producing a distribution of outcomes rather than a single equilibrium price.  The observed persistent dispersion should therefore be understood as positive structural validation: the simulator faithfully instantiates the expected property of decentralised bilateral bargaining, rather than inadvertently collapsing to a degenerate uniform outcome.

\paragraph{Finding 7: Liquidity is high.}
Mean liquidity of 0.81 (81\%) indicates that the large majority of matched pairs reach agreement, suggesting that the Qwen2.5-14B model is effective at extracting deals within 10 rounds under the given parameter distributions without gate assistance.

\paragraph{Finding 8: Surplus is approximately symmetric.}
Mean buyer surplus and seller surplus are approximately balanced (\$31.48 and \$31.68 respectively), consistent with the scenario parameters (buyer value from [\$80, \$150], seller cost from [\$50, \$120], creating varied ZOPA distributions).  H-3c is not supported: the first-mover advantage observed in fixed-scenario sessions (Experiment~1) does not translate into asymmetric aggregate surplus in the market setting with heterogeneous parameters.  This finding is in direct contrast to the 5.65:1 buyer-to-seller concession ratio observed in Experiment~1, and the contrast is itself informative.  A plausible explanation is that under heterogeneous parameter draws, the randomness of the value and cost distributions dominates any systematic role-based asymmetry: in some pairings the buyer holds most of the bargaining leverage, in others the seller does, and these effects wash out in the aggregate.  The buyer-first concession advantage observed under fixed-scenario conditions is therefore a property of that specific parameterisation, not a robust feature of the protocol when agents negotiate under varied private constraints.

% ───────────────────────────────────────────────────────────────────
\subsubsection{Experiment 4: Stochastic Shock Response}
\label{sec:results-shock}

Experiment~4 applies stochastic multiplicative shocks (30\% per-tick probability; demand and supply multipliers drawn uniformly from [0.7, 1.3]) to a 10-tick market and compares outcomes against an unshocked baseline.  600 sessions were run (300 per condition, 3 seeds).  Deal rate is 70.5\% (423/600).

\begin{table}[ht]
\centering
\caption{Shock response: difference-in-differences (DiD) analysis (300 sessions per condition, 3 seeds).  Pre-shock period: ticks~0--4; post-shock period: ticks~5--9.  DiD is a structured descriptive comparison, not a causal estimate.}
\label{tab:shock}
\begin{tabular}{@{}lcccc@{}}
\toprule
Condition & Pre-period mean price (\$) & Post-period mean price (\$) & Within-condition diff (\$) & DiD (\$) \\
\midrule
\texttt{no\_shock}   & 107.88 & 110.62 & $+$2.75 & \\
\texttt{with\_shock} & 111.44 & 107.66 & $-$3.78 & $-$6.53 \\
\bottomrule
\end{tabular}
\end{table}

Liquidity DiD: $-0.060$ (marginal suppression of deal rates under shocks).

\paragraph{Finding 9: Shocks produce a price-depressing effect of approximately \$6.53.}
The structured descriptive DiD comparison yields $-$\$6.53: the shock condition declines by \$3.78 across the two halves while the no-shock condition rises by \$2.75.  This is directionally consistent with the prediction that mixed-direction stochastic shocks depress average prices and reduce liquidity, since negative shocks (demand multiplier $<$1 or supply multiplier $>$1) compress the ZOPA and reduce the floor available for deal completion.  H-4b is supported directionally.

\paragraph{Finding 10: No consistent price recovery within the 10-tick window.}
Seed-level trajectories show heterogeneous behaviour: seed~42 exhibits partial recovery after the shock period, seed~123 shows continued decline, and seed~456 shows mild depression without recovery.  There is no consistent recovery pattern.  This finding should be treated as suggestive rather than conclusive: with only three seeds and 10 ticks each, the seed-to-seed variance dominates the aggregate pattern.

\paragraph{Finding 11: Liquidity impact is marginal.}
The DiD for liquidity is $-0.060$, a modest reduction compared to the price effect.  Stochastic shocks of symmetric magnitude (multipliers in [0.7, 1.3] with mean 1.0) produce only small net liquidity effects on average, since positive and negative shocks partially offset.  Tail events (large negative shocks) will produce locally severe liquidity drops within individual ticks, but these are diluted in the aggregate.  H-4a is supported directionally: the shock condition shows greater tick-to-tick price variance than baseline.  H-4c (asymmetric tail risk) cannot be evaluated with only three seeds, as within-seed variance dominates any tail analysis.

\paragraph{Interpretive caveat.}
The DiD comparison is a structured descriptive analysis.  It is not a causal estimate: the pre- and post-periods coincide with a time trend in the no-shock condition (+\$2.75), and the shock condition's evolution is influenced by the specific shock realisations drawn under each seed.  The \$6.53 estimate should be understood as a directional signal under these particular parameter settings, not as a precise measure of shock impact.  A post-hoc pooled comparison (Section~\ref{sec:results-anchor}) reports a different quantity---the mean price difference between shock and no-shock sessions pooled across all ticks---and is not contradictory to this DiD estimate.

% ───────────────────────────────────────────────────────────────────
\subsubsection{Experiment 4 Extension: Anchor Environment}
\label{sec:results-anchor}

The Experiment~4 anchor-environment extension tests whether the item reference price in the agent prompt functions as an environmental price anchor that shapes both price levels and shock responsiveness.  Three anchor modes---\texttt{fixed\_anchor} (constant reference price, the default from Experiment~4), \texttt{updated\_anchor} (dynamically updated to the prior tick's mean deal price), and \texttt{no\_anchor} (reference price omitted)---are each crossed with the shock/no-shock condition from Experiment~4.  The \texttt{fixed\_anchor} condition reuses the Experiment~4 dataset; the other two modes add 1{,}200 new sessions.

\begin{table}[ht]
\centering
\caption{Experiment~4 extension: informational anchor environment results.  Mean price and liquidity are reported for the no-shock condition; shock price effect is the difference in mean deal price between the with-shock and no-shock conditions within each anchor mode.  Fixed-anchor values are derived from the Experiment~4 dataset.}
\label{tab:anchor-environment}
\begin{tabular}{@{}lcccc@{}}
\toprule
Anchor mode & Mean price, no-shock (\$) & Price std (\$) & Liquidity & Shock price effect (\$) \\
\midrule
\texttt{fixed\_anchor}   & $\approx$109 & $\approx$15  & 0.69 & $-$6.53 \\
\texttt{updated\_anchor} & 78.84    & 14.16 & 0.713 & $+$4.15 \\
\texttt{no\_anchor}      & 99.21    & 7.73  & 0.867 & $-$0.20 \\
\bottomrule
\end{tabular}
\end{table}

\paragraph{Finding 12: The reference price signal is a primary determinant of price levels.}
Across the three anchor modes, baseline mean prices span from approximately \$79 to \$109---a range of roughly \$30---despite identical buyer value and seller cost distributions in all conditions.  H-4ext-a is strongly supported: the environmental reference signal shapes price levels substantially, exceeding the \$4.63 price response to a \$20 demand shock reported in Experiment~6 (Section~\ref{sec:results-supply-demand}).  This finding demonstrates that LLM agents are highly sensitive to the informational environment in which they negotiate, not only to their private economic constraints.  The \$30 spread is a strong result: it means that under these experimental conditions, the informational environment explains more price variance than the underlying economic fundamentals do.  This has a direct methodological implication: any LLM-based market simulation that embeds a reference price in the agent prompt is, intentionally or not, pre-determining a substantial portion of the price level.  The reference signal is not a neutral contextual detail---it is a first-order design variable.  This finding also helps explain the price stickiness observed in Experiment~6.  If agents anchor to a prompt-embedded reference price rather than independently reasoning from their private constraints, then a \$20 shift in those constraints will produce only a muted price response: the agent continues to negotiate toward the anchored reference level rather than fully internalising the changed fundamental.  The low elasticities in Experiment~6 (0.275 and 0.035) are thus plausibly a downstream consequence of the reference-price anchoring mechanism identified here.

\paragraph{Finding 13: Updated anchors create self-reinforcing price drift.}
Under \texttt{updated\_anchor}, baseline mean prices settle at \$78.84, well below both the \texttt{fixed\_anchor} ($\sim$\$109) and \texttt{no\_anchor} (\$99.21) conditions.  The mechanism is a self-reinforcing feedback loop: below-equilibrium deal prices in early ticks feed a lower reference price into subsequent prompts, which in turn produces lower opening offers and sustains further below-mean settlements.  The shock price effect under this mode is \emph{positive} (+\$4.15), reversing the sign observed in the fixed-anchor condition ($-$\$6.53).  A plausible behavioural interpretation is that stochastic shocks partially disrupt the downward feedback mechanism, producing temporarily higher reference prices that lift the overall mean.  The sign reversal cautions against treating the direction of the shock effect as a robust property of LLM negotiation behaviour independent of the informational anchor environment.

\paragraph{Finding 14: Removing the anchor stabilises prices but eliminates shock sensitivity.}
The \texttt{no\_anchor} condition produces the most stable price series (std~=~\$7.73 vs.\ \$14.16 under updated-anchor) and the highest liquidity (0.867 vs.\ 0.69--0.71).  However, the shock price effect is effectively zero ($-$\$0.20), indicating that without a reference signal, agents default to an internal prior and do not update it in response to fundamental changes.  H-4ext-b is supported: removing the anchor eliminates shock sensitivity.

\paragraph{Supplementary statistical check: pooled anchor-mode comparison.}
The three anchor-mode conditions were compared using a pooled post-hoc analysis of 1{,}388 completed deals drawn from Experiment~4 (fixed-anchor arm, both shock conditions) and the Experiment~4 extension (updated-anchor and no-anchor arms, both shock conditions) across three seeds.  Fixed-anchor data originate from Experiment~4; updated-anchor and no-anchor data from the Experiment~4 extension.  The analysis uses Welch's $t$-tests with Bonferroni correction for three pairwise comparisons (corrected $\alpha = 0.0167$).  This is a post-hoc descriptive analysis, not a pre-registered confirmatory test.  An important methodological caveat applies: sessions within a seed share market context, violating the independence assumption underlying the standard error estimates; consequently, reported standard errors may be underestimated and $p$-values should be interpreted with commensurate caution.

Under the tested model--prompt--protocol configuration, all three pairwise price-level comparisons are statistically significant after Bonferroni correction (Table~\ref{tab:anchor-pairwise}).  Effect sizes range from medium to large by conventional benchmarks ($d = 0.48$ to $d = 1.32$), consistent with the qualitative price-level ordering observed in Table~\ref{tab:anchor-environment}.

\begin{table}[ht]
\centering
\caption{Post-hoc pairwise comparison of mean transaction prices across anchor modes (1{,}388 deals pooled over shock conditions, Welch's $t$-test with Bonferroni-corrected $\alpha = 0.0167$).  Fixed-anchor data from Experiment~4; updated-anchor and no-anchor data from the Experiment~4 extension.}
\label{tab:anchor-pairwise}
\begin{tabular}{@{}lcccc@{}}
\toprule
Comparison & $\Delta$ mean (\$) & 95\% CI (\$) & $p$ & Cohen's $d$ \\
\midrule
fixed vs.\ updated    & $+$26.79 & [$+$24.05, $+$29.53] & $1.5 \times 10^{-68}$ & $+$1.32 \\
fixed vs.\ no\_anchor  & $+$10.01 & [$+$7.35, $+$12.67]  & $3.6 \times 10^{-13}$ & $+$0.48 \\
updated vs.\ no\_anchor & $-$16.78 & [$-$19.33, $-$14.23] & $3.8 \times 10^{-35}$ & $-$0.83 \\
\bottomrule
\end{tabular}
\end{table}

Within-anchor-mode shock sensitivity was assessed by comparing pooled mean transaction prices between the with-shock and no-shock sessions for each mode.  The observed differences are: fixed-anchor $+$\$0.60 ($p = 0.77$, n.s.), updated-anchor $+$\$4.07 ($p = 0.032$, $d = 0.21$), and no-anchor $+$\$0.18 ($p = 0.92$, n.s.).  Only the updated-anchor condition shows a statistically significant mean price difference between shock conditions, though the effect size is small ($d = 0.21$) and should be treated with the same non-independence caveat noted above.

A reconciliation with Finding~9 (Section~\ref{sec:results-shock}) is warranted.  The pooled fixed-anchor shock effect of $+$\$0.60 (n.s.) does not contradict the DiD estimate of $-$\$6.53 reported there.  The DiD measures temporal price \emph{change} across a pre/post split (ticks~0--4 vs.\ 5--9), capturing within-condition dynamics over time.  The pooled shock comparison reported here measures the overall mean price \emph{level} difference between shock and no-shock sessions across all ticks, collapsing the temporal dimension.  Both statistics are valid descriptions of different aspects of the same data.  Together, they suggest that the fixed-anchor shock effect manifests as a directional temporal shift within each run rather than as a persistent mean price elevation detectable in the session-level pooled comparison.

\paragraph{Behavioural interpretation.}
The \$30 price-level spread across anchor modes, combined with the sign reversal in the shock effect, indicates that LLM price-setting behaviour is strongly conditioned on the reference information presented in the prompt.  Agents do not independently discover prices from private constraints alone; they use the environmental reference signal as an anchor around which to negotiate.  This behaviour is consistent with the anchoring literature~\cite{tversky1974,galinsky2001}, though the mechanism differs: in human studies, anchoring operates through an opponent's first offer, whereas here it operates through a prompt-embedded environmental signal.  For experimental design, this implies that anchor mode must be explicitly controlled and reported as a first-order design variable.

% ───────────────────────────────────────────────────────────────────
\subsubsection{Experiment 6: Structural Supply-Demand Shifts}
\label{sec:results-supply-demand}

Experiment~6 applies deterministic, directional \$20 shifts to buyer values (demand shock) or seller costs (supply shock) and compares outcomes against a baseline.  576 sessions were run (192 per condition, 3 seeds).  Deal rate is 70.5\% (406/576).

\paragraph{Welfare accounting identity caveat.}
Before reporting welfare figures, a structural caution is necessary.  Total welfare in this simulator is defined as the sum of buyer surplus and seller surplus: $(v_b - p) + (p - c_s) = v_b - c_s$, where $v_b$ is buyer value, $c_s$ is seller cost, and $p$ is the deal price.  The price $p$ cancels.  Consequently, aggregate welfare is mechanically determined by the distribution of $(v_b, c_s)$ pairs that reach deals, not by the price at which those deals close.  Welfare comparisons across supply-demand conditions reflect changes in the composition of trading pairs, not agent-level bargaining behaviour.  This identity must be borne in mind when evaluating the findings below.  Price movements, by contrast, require genuine LLM behavioural adaptation and are therefore the more informative signal of agent price discovery.

\begin{table}[ht]
\centering
\caption{Supply-demand experiment outcomes by condition (192 sessions per condition, 3 seeds).  Buyer value and seller cost columns show mean values across the distribution used in each condition.}
\label{tab:supply-demand}
\begin{tabular}{@{}lcccccc@{}}
\toprule
Condition & Mean price (\$) & Liquidity & Total welfare (\$) & Buyer surplus (\$) & Seller surplus (\$) & Buyer share (\%) \\
\midrule
\texttt{baseline}      & 109.33 & 0.724 & 65.60 & 21.40 & 44.19 & 33 \\
\texttt{demand\_shock} & 113.96 & 0.724 & 87.24 & 39.07 & 48.18 & 45 \\
\texttt{supply\_shock} & 110.50 & 0.667 & 49.84 & 22.16 & 27.67 & 44 \\
\bottomrule
\end{tabular}
\end{table}

\begin{table}[ht]
\centering
\caption{Price elasticity of supply and demand shocks.  Elasticity is defined as proportional price change divided by proportional fundamental change (a $\pm$\$20 shift on a mean base).}
\label{tab:elasticities}
\begin{tabular}{@{}lcccc@{}}
\toprule
Shock & Input shift (\$) & Price change (\$) & Price change (\%) & Elasticity \\
\midrule
Demand shock & $+$20 (buyer value) & $+$4.63 & $+$4.2\% & 0.275 \\
Supply shock & $+$20 (seller cost) & $+$1.17 & $+$1.1\% & 0.035 \\
\bottomrule
\end{tabular}
\end{table}

\paragraph{Finding 15: Prices exhibit strong stickiness under both shocks.}
Despite \$20 shifts in underlying fundamentals, transaction prices move by only \$4.63 under the demand shock (+4.2\%) and \$1.17 under the supply shock (+1.1\%).  The implied price elasticities are 0.275 and 0.035 respectively---far below the full pass-through (elasticity of 1.0) that classical price theory would predict.  Agents appear to anchor to a stable reference price range and adjust it only weakly in response to changing fundamentals, even when those fundamentals are encoded in their private constraints.  This price stickiness is the most substantively important finding in Experiment~6, because it requires the LLM to update prices in response to changed fundamentals rather than following mechanically from the accounting identity.

\paragraph{Finding 16: Demand shock raises prices; supply shock barely affects them.}
The demand shock (+\$20 buyer value) produces a \$4.63 price increase, four times larger than the supply shock's \$1.17 despite identical input magnitudes.  H-6a is supported directionally: the demand shift raises mean price, though the magnitude is far below full pass-through.  H-6b is partially supported: the supply shift raises prices but by a negligible amount (\$1.17), and liquidity drops from 0.724 to 0.667, consistent with ZOPA compression.  This asymmetry is directionally inconsistent with classical competitive market predictions, in which a cost-floor increase (supply shock) should produce stronger price pass-through than a value increase (demand shock), because sellers can enforce a hard minimum floor at their cost.  One possible explanation is that buyers drive pricing more aggressively than sellers in this bargaining protocol---consistent with the buyer's first-mover advantage---so that demand-side changes propagate into prices more readily.  An alternative explanation is model-specific anchoring to cost-adjacent prices.  Neither explanation is resolvable without targeted ablation.

\paragraph{Finding 17: Welfare changes track the accounting identity.}
Total welfare rises by \$21.65 under the demand shock (\$65.60 to \$87.24) and falls by \$15.76 under the supply shock (\$65.60 to \$49.84).  As noted in the caveat above, these changes are mechanically determined by the shift in $(v_b - c_s)$ distributions and are not evidence of agent-level welfare optimisation.  Liquidity is unchanged under the demand shock (0.724 in both conditions), indicating that the wider value distribution does not translate into more deals---a further sign of price stickiness, because buyers with higher values could afford to accept the seller's price but do not consistently close faster.

\paragraph{Finding 18: Supply shock redistributes surplus toward buyers, relative to total.}
Under the supply shock, buyer surplus rises marginally (+\$0.76, from \$21.40 to \$22.16) while seller surplus falls sharply ($-$\$16.52, from \$44.19 to \$27.67).  The buyer share of total welfare increases from 33\% to 44\%.  This distributional shift is consistent with cost-floor pressure: higher seller costs reduce the available seller margin at the prices at which deals are closing, even though the closing price itself barely changes.  Sellers absorb the cost shock with little ability to pass it forward through higher prices.  H-6c is directionally supported: the liquidity drop under the supply shock (0.724 to 0.667) exceeds what ZOPA compression alone would predict, suggesting a behavioural component in which agents fail to reach agreement even when a narrow positive ZOPA exists.

% ───────────────────────────────────────────────────────────────────
\subsubsection{Experiment 5: Matching Mechanism Comparison}
\label{sec:results-mechanism}

Experiment~5 compares random matching (baseline) against surplus-maximising (oracle) matching in an 8-tick market.  335 sessions were run (random: 192; surplus-max: 143; the lower surplus-max count arises because the matcher excludes negative-ZOPA pairs).  Overall deal rate is 83.9\%.

\begin{table}[ht]
\centering
\caption{Mechanism comparison outcomes (3 seeds, 8 ticks per seed).  Total surplus is the sum of buyer and seller surplus across all deals.}
\label{tab:mechanism}
\begin{tabular}{@{}lccc@{}}
\toprule
Metric & \texttt{random} & \texttt{surplus\_max} & Difference \\
\midrule
Total sessions           & 192     & 143     & $-$49 \\
Total surplus (\$)       & 49.48   & 53.05   & $+$3.57 ($+$7.2\%) \\
Price std (\$)           & 17.92   & 16.14   & $-$1.78 \\
Liquidity                & 0.85    & 0.82    & $-$0.03 \\
Mean buyer surplus (\$)  & 25.52   & 26.61   & $+$1.09 \\
Mean seller surplus (\$) & 23.96   & 26.43   & $+$2.47 \\
\bottomrule
\end{tabular}
\end{table}

\paragraph{Finding 19: Surplus-maximising matching improves total welfare by 7.2\%.}
Total surplus rises from \$49.48 under random matching to \$53.05 under surplus-max matching, a 7.2\% improvement.  The surplus-max matcher achieves this by directing trades exclusively toward pairs with positive ZOPA overlap, increasing the probability that any given session produces a deal.  H-5b is supported.

\paragraph{Finding 20: Mechanism design benefits both parties.}
Both buyer surplus (+\$1.09) and seller surplus (+\$2.47) increase under the improved mechanism.  Seller surplus gains are proportionally larger.  This pattern is consistent with the scenario producing a balanced surplus environment, where sellers can negotiate outcomes closer to their preferred range.

\paragraph{Finding 21: Liquidity is marginally lower under surplus-max, yet welfare is higher.}
Liquidity falls slightly from 0.85 to 0.82 despite the mechanism filtering out low-ZOPA pairs.  This counter-intuitive result may arise because surplus-max occasionally creates highly asymmetric pairings in which one side holds all the bargaining leverage, making agreement harder to reach even when ZOPA overlap exists.  The net effect is positive for welfare because the deals that do close are of higher quality.  H-5a (higher deal rate under surplus-max) is not supported; the mechanism improves welfare quality rather than quantity.

% ═══════════════════════════════════════════════════════════════════
%  Cross-Experiment Synthesis
% ═══════════════════════════════════════════════════════════════════

\subsection{Cross-Experiment Synthesis}
\label{sec:results-cross}

This section synthesises patterns across experiments and organises findings by evidence strength.

\paragraph{Theme 1: Asymmetric concession burden as a protocol-level pattern.}
Across Experiments~1, 2, and~6, a consistent pattern emerges: buyers concede more than sellers.  In Experiment~1, the buyer-to-seller concession ratio is 5.65:1, and all deals cluster in the bottom 28\% of the ZOPA.  In Experiment~2, buyer surplus exceeds 83\% of the ZOPA across all deadline conditions, with the seller barely moving from its initial position.  In Experiment~6, buyer surplus is lower than seller surplus at baseline (33\% vs.\ 67\% of total welfare), and the supply shock redistributes only modestly toward buyers.  This asymmetry is consistent with the buyer-first turn order in the alternating-offer protocol and with the observed seller anchoring behaviour, but it is also sensitive to the specific parameter settings used.  Whether this asymmetry is an intrinsic property of LLM bargaining or an artefact of the scenario configuration cannot be determined without ablation of the role-assignment and parameter settings.

\paragraph{Theme 2: Welfare as accounting identity versus price as behavioural signal.}
A recurring interpretive challenge across Experiments~6 and~5 is that welfare measures partly reflect the accounting identity $(v_b - c_s)$ rather than agent decision quality.  Welfare increases when the distribution of $(v_b, c_s)$ pairs shifts favourably, and when trades are concentrated in high-surplus pairs (as in surplus-max matching), regardless of how agents negotiate.  Price movements, by contrast, require the LLM agent to adjust its offer in response to changed fundamentals.  The consistently low price elasticities in Experiment~6 (0.275 for demand shocks, 0.035 for supply shocks) and the persistent within-tick price dispersion in Experiment~3 indicate that price discovery is weak: agents do not efficiently update prices in response to shifting values and costs.  This is the more diagnostically useful finding about LLM bargaining capability.  This distinction---between welfare metrics that track an accounting identity and price metrics that require genuine agent adaptation---should inform how future LLM negotiation studies evaluate market outcomes: reporting welfare without separately reporting price elasticity risks attributing to agent capability what is in fact a structural property of the parameter design.

\begin{table}[ht]
\centering
\caption{Cross-experiment summary of findings, organised by evidence strength.  ``Strong'' indicates a finding consistent across all three seeds with directional clarity; ``Moderate'' indicates a finding with directional support but high cross-seed variance or limited replication; ``Weak'' indicates a finding from a pattern that is unclear or not replicated.}
\label{tab:synthesis}
\small
\begin{tabular}{@{}llcl@{}}
\toprule
Pattern & Experiments & Strength & Primary driver \\
\midrule
Concave (Boulware) concession curves & 1 & Strong & Agent \\
Asymmetric concession burden (5.65:1) & 1, 2 & Strong & Protocol/Agent \\
Deadline acceleration (82\% last-2 at 6 rounds) & 2 & Strong & Agent \\
Reference-signal price sensitivity ($\sim$\$30 spread) & 4ext & Strong & Prompt/Design \\
Surplus-max mechanism $\to$ +7.2\% welfare & 5 & Moderate & Mechanism \\
Shock price depression (DiD $-$6.53) & 4 & Moderate & Protocol/Agent \\
Anchor-mode moderation of shock sensitivity & 4, 4ext & Moderate & Prompt/Design \\
Price stickiness (elasticity 0.03--0.28) & 6 & Moderate & Agent \\
Welfare tracks fundamentals (accounting identity) & 6 & (Mechanical) & Design \\
Price convergence in market & 3 & Weak & None observed \\
Consistent shock recovery & 4 & Weak & Not observed \\
\bottomrule
\end{tabular}
\end{table}

\paragraph{Theme 3: Deal success rates vary systematically across experimental conditions.}
Deal success rates range from 97.2\% in fixed-scenario experiments (Experiment~1) to 70.5\% in distribution-mode market experiments (Experiments~4 and~6).  Table~\ref{tab:deal-rates} summarises deal rates across all experiments.

\begin{table}[ht]
\centering
\caption{Deal success rates across all experiments.}
\label{tab:deal-rates}
\small
\begin{tabular}{@{}llcc@{}}
\toprule
Exp & Name & Deal rate (\%) & Sessions \\
\midrule
1 & Concession (LLM) & 97.2 & 72 \\
1 & Concession (rule-based) & 97.2 & 72 \\
2 & Deadline (6 rounds) & 94.4 & 72 \\
2 & Deadline (12 rounds) & 97.2 & 72 \\
2 & Deadline (20 rounds) & 97.2 & 72 \\
3 & Market dynamics & 82.3 & 300 \\
4 & Shock (no shock) & $\approx$69 & 300 \\
4 & Shock (with shock) & $\approx$72 & 300 \\
6 & Supply-demand (baseline) & 72.4 & 192 \\
6 & Supply-demand (demand shock) & 72.4 & 192 \\
6 & Supply-demand (supply shock) & 66.7 & 192 \\
4ext & Anchor (fixed, no shock) & $\approx$69 & 300 \\
4ext & Anchor (updated, no shock) & 71.3 & 300 \\
4ext & Anchor (no anchor, no shock) & 86.7 & 300 \\
5 & Mechanism (random) & 85.0 & 192 \\
5 & Mechanism (surplus-max) & 82.0 & 143 \\
\bottomrule
\end{tabular}
\end{table}

The primary determinant of deal success is the experimental setting: fixed-scenario experiments with wide ZOPAs (Experiments~1 and~2) achieve near-complete agreement, while distribution-mode experiments produce lower rates because some randomly matched pairs have narrow or absent ZOPAs.  Among market experiments, the \texttt{no\_anchor} condition in the Experiment~4 extension achieves the highest deal rate (86.7\%), suggesting that removing the reference price signal may reduce anchoring-induced rigidity that prevents agreement.  The supply shock in Experiment~6 produces the lowest deal rate (66.7\%), consistent with ZOPA compression from higher seller costs.

\paragraph{Theme 4: Comparison to human negotiation patterns.}
Several observed behavioural patterns are directionally consistent with patterns documented in the human negotiation literature, though all comparisons are indirect (no human participants were tested under the same protocol).

\emph{Deadline-driven concession acceleration.}  The Experiment~2 finding---82.4\% of deals close in the final two rounds under a 6-round horizon, compared to 11.4\% under 20 rounds---is directionally consistent with Roth et al.'s~\cite{roth1988deadline} documentation of deadline effects in human bargaining.  However, LLM agents do not adjust final prices across deadline conditions (\$84.93 to \$86.73, a \$2 range), whereas human deadline studies typically document both timing and price effects.  This dissociation---timing sensitivity without price sensitivity---is a specific, testable behavioural difference between LLM and human negotiators.  It suggests that the LLM responds to deadline salience primarily as a coordination signal (``the window is closing; close the deal now'') rather than as a pressure signal that alters its assessment of an acceptable price.  Whether this difference is model-specific, prompt-specific, or a general property of LLM bargaining is a question that could be probed directly through studies that vary both the deadline condition and the human baseline under the same protocol.

\emph{Boulware-style concession.}  The concave concession trajectory observed in Experiment~1 (large initial concession followed by rapid hardening) is consistent with the Boulware bargaining pattern documented in the negotiation literature~\cite{raiffa1982}.  The prompt instructs monotonic concession direction but does not dictate step size or timing; the frontloaded concession pattern is therefore partly attributable to LLM behaviour, though the boundary between prompt-induced and agent-driven dynamics cannot be determined from current experiments.

\emph{Anchor sensitivity.}  The Experiment~4 extension demonstrates that LLM agents are strongly conditioned on the environmental reference price, consistent with the anchoring literature~\cite{tversky1974,galinsky2001}, though the mechanism differs: in human studies, anchoring operates through an opponent's first offer, whereas here it operates through a prompt-embedded environmental signal.

\paragraph{Protocol-level considerations.}
Across all experiments, the bargaining protocol is held constant: agents operate under a finite-horizon counteroffer protocol with structured action extraction, monotonic concession instructions, and pre-computed feasibility information.  These protocol choices constrain concession direction and steer toward deal finalization when offers are feasible.  The consistently low seller concession observed in Experiment~1 may partly reflect prompt asymmetry.  Prompt-level ablation studies---testing deal rates and concession dynamics with and without the monotonic concession instruction, the accept-condition block, and the reference price signal---would further clarify the boundary between agent-driven and protocol-induced behaviour.
